\section{Related Work}
\label{sec:related}

\paragraph{Do-All and asynchronous PRAMs.} Our method belongs with the work on Write-All, Do-All and Certified Write-All~\cite{kanellakis_robust, georgiou_doall, martel_cwa} (Section~\ref{sec:bg-doall}). Some algorithms for these problems let processors walk the tasks in permutations of their own and record their progress in shared memory, as our threads walk buckets in random orders and publish their progress and certificates, and they bound the work that processors repeat by the contention of the permutations~\cite{anderson_woll_cwa}, which our analysis adapts (Section~\ref{sec:pf-work}). A CWA per step runs any PRAM program on asynchronous processors~\cite{martel_apram, kanellakis_book}, so these results give a wait-free BFS, most simply with one CWA over all $n$ vertices per level, at work $\Omega(Dn)$, a bound of this mapping rather than of every construction built on CWA; to our knowledge, none of these constructions was implemented to compete with tuned parallel BFS. Our method differs from them in that its tasks are discovered during the computation, by the same threads, and grouped by the thread that discovered them, so that a thread certifies a level by reading $T$ flags rather than $n$ cells; in that its tasks are weighted by degrees, so that what the others redo for a stopped thread depends on how finely the work is certified (Proposition~\ref{prop:granularity}); and in that its threads agree on a plan for each level with one CAS. Table~\ref{tab:related} compares these approaches: among the wait-free ones, ours is the only one whose threads share the work of a top-down level at a cost that depends on the level rather than on the graph.

\paragraph{Wait-free objects and tasks.} Wait-freedom is studied mostly where inputs are distributed, for concurrent objects~\cite{herlihy_waitfree, aomp, kogan_fastwf} and for decision tasks such as consensus~\cite{flp} and set agreement, whose wait-free solvability Herlihy and Shavit characterize~\cite{herlihy_shavit_topology}; in a search, it is a question of work, not of solvability (Section~\ref{sec:bg-doall}). Wait-freedom has also been studied for concurrent graph data structures. Kallimanis and Kanellou~\cite{kallimanis_graph} provide wait-free graph objects with consistent dynamic traversals, using helping to complete pending operations, and Bhardwaj et~al.~\cite{graph_snapshots} use multiversioning to take wait-free snapshots of a lock-free dynamic graph, on which queries such as BFS can run during concurrent updates. These works address consistent operations and traversals of a changing graph, each run by one thread, while our threads cooperate on a single search of a graph that does not change.

\paragraph{Parallel BFS\@.} The parallel BFS methods for shared memory are mostly level-synchronous, from the fine-grained multithreaded BFS of Bader and Madduri for the Cray MTA-2~\cite{bader_bfs} and the work-efficient BFS of Leiserson and Schardl, which builds frontiers with bags and Cilk++ reducers~\cite{charles_bfs}, to the direction-optimizing BFS of Beamer et~al.~\cite{beamer_dobfs}, which GAPBS implements~\cite{gapbs} and whose direction rule our method uses, and the vertex subsets of Ligra~\cite{ligra} and GBBS~\cite{gbbs}, built on Ligra. PASGAL~\cite{pasgal} addresses the synchronization and scheduling overhead of large-diameter graphs with the vertical granularity control and hash bags of Wang et~al.~\cite{vgc_hashbag}: a task explores several hops locally, which reduces the number of global rounds, and vertices may be revisited when shorter distances are found. All of these are blocking: they end every step at a barrier or a join.

\paragraph{BFS without barriers.} Other parallel BFS methods relax the order between levels. Pearce et~al.~\cite{pearce_async} traverse graphs asynchronously in shared memory, for BFS, shortest paths from a source, and connected components: threads process visits to vertices from queues, with no barriers between levels, which hides the latency of memory and of solid-state storage, and visit a vertex again when they find a shorter path to it. On GPUs, XBFS~\cite{xbfs} makes the bottom-up step asynchronous, so that one iteration can discover vertices of more than one level, but each iteration still starts once the one before has ended. Unordered BFS gives up levels altogether, at the cost of potentially more work; Holst and Selin~\cite{holst_unordered} schedule it with a relaxed FIFO queue, and report speedups of 1.6 to 2.4 over a level-synchronous BFS on sparse, high-diameter graphs. None of these makes a search wait-free: XBFS waits for its slowest thread at the end of every iteration, and in an asynchronous or unordered BFS without helping, a thread that is suspended after taking a vertex or a visit from a queue holds work that no other thread will do, so that the search cannot complete until it resumes.

\paragraph{Work stealing.} Work stealing~\cite{blumofe_workstealing} balances load by letting a thread take tasks from the queues of others, and hands out every task exactly once, so a thread that stops after taking a task holds it for good. Idempotent work stealing~\cite{michael_iws} hands a task out at least once, which lets the owner's operations avoid atomic instructions and store-load fences, for applications that tolerate or detect repeated tasks, such as transitive closure and spanning trees. Our methods tolerate repeated work as well, but that alone does not make a search wait-free: a task that a stopped thread has taken may be handed out again, but nothing ensures that it will be. What our methods add is an order of publishing work and progress in which a stopped thread leaves nothing that the others cannot redo: the work stays in its bucket or block, where the others redo what has not been certified, and each thread decides alone when a level is complete.
